\section{Especificación léxica}

Se define $L=[a\text{-}z]$, $U=[A\text{-}Z]$, $D=[0\text{-}9]$ y $C=L\cup U\cup D\cup\{\_\}$. El alfabeto de entrada $\Sigma$ contiene los caracteres Unicode de la fuente, pero los identificadores, dígitos y espacios admitidos se restringen a las clases ASCII indicadas. Los caracteres Unicode adicionales pueden aparecer dentro de citas y comentarios. LF, CRLF y CR se cuentan como un único salto de línea; una tabulación ocupa una columna. Las posiciones comienzan en uno y cuentan caracteres, no bytes ni anchura visual.

Las expresiones de la Tabla~\ref{tab:formal} usan la notación de \texttt{re} de Python: $|$ denota alternativas, los corchetes definen clases y los grupos \texttt{(?:...)} no capturan. Los patrones se aplican desde la posición actual mediante \texttt{match}; no se busca una coincidencia posterior que omita caracteres. En Python las alternativas se intentan de izquierda a derecha, de modo que el orden de operadores largos antes de sus prefijos forma parte de la implementación de máxima coincidencia \cite{pythonre}. Las citas y los comentarios de bloque se recorren explícitamente para conservar el fragmento de los errores de cierre.

\small
\begin{longtable}{L{3.1cm}L{4.7cm}L{7.1cm}}
\caption{Definición formal del subconjunto.}\label{tab:formal}\\
\toprule
\tableheader{Token} & \tableheader{Descripción} & \tableheader{Expresión regular conceptual} \\
\midrule
\endfirsthead
\toprule
\tableheader{Token} & \tableheader{Descripción} & \tableheader{Expresión regular conceptual} \\
\midrule
\endhead
ATOMO & Identificador iniciado en minúscula; se reclasifican \texttt{is} y \texttt{mod}. & \path{[a-z][A-Za-z0-9_]*} \\
\rowcolor{palegray} \makecell[l]{ATOMO\_\\CITADO} & Texto entre comillas simples con escapes admitidos. & \path{'(?:\\[ntr\\'"]|[^'\\\r\n])*'} \\
VARIABLE & Identificador iniciado en mayúscula o guion bajo, excepto el lexema aislado \texttt{\_}. & \path{(?:[A-Z][A-Za-z0-9_]*|_[A-Za-z0-9_]+)} \\
\rowcolor{palegray} \makecell[l]{VARIABLE\_\\ANONIMA} & Guion bajo aislado. & \path{_} \\
ENTERO & Secuencia decimal sin signo. & \path{[0-9]+} \\
\rowcolor{palegray} REAL & Decimal o exponente; el signo solo se admite en el exponente. & \path{(?:[0-9]+\.[0-9]+(?:[eE][+-]?[0-9]+)?|[0-9]+[eE][+-]?[0-9]+)} \\
CADENA & Texto entre comillas dobles con escapes admitidos. & \path{"(?:\\[ntr\\'"]|[^"\\\r\n])*"} \\
\rowcolor{palegray} \makecell[l]{OPERADOR\_\\CLAUSULA} & Operadores de regla y gramática de cláusulas definidas. & \path{-->|:-} \\
\makecell[l]{OPERADOR\_\\CONSULTA} & Inicio de consulta. & \path{\?-} \\
\rowcolor{palegray} \makecell[l]{OPERADOR\_\\COMPARACION} & Unificación, identidad, desigualdad y comparación. & \path{\\==|=\.\.|\\=|==|=<|>=|=|<|>} \\
\makecell[l]{OPERADOR\_\\ARITMETICO} & Operadores simbólicos y alfabéticos admitidos. & \path{//|\*\*|[+\-*/]|is|mod} \\
\rowcolor{palegray} \makecell[l]{OPERADOR\_\\CONTROL} & Negación, corte y disyunción. & \path{\\\+|!|;} \\
DELIMITADORES & Familia de tokens detallada en la Tabla~\ref{tab:delimitadores}. & \path{[()\[\]{}|,.]} \\
\rowcolor{palegray} COMENTARIO de línea & Desde \texttt{\%} hasta antes de CR, LF o fin de archivo; se descarta. & \path{%[^\r\n]*} \\
COMENTARIO de bloque & Hasta el primer cierre; sin anidamiento; se descarta. & \path{/\*[^*]*\*+(?:[^/*][^*]*\*+)*/} \\
\rowcolor{palegray} ESPACIO & Separadores descartados con actualización de posición. & \path{[ \t\r\n\f\v]+} \\
\bottomrule
\end{longtable}
\normalsize

Los operadores cubren los lexemas obligatorios de la tarea y la extensión opcional \texttt{-->}. Prolog permite declarar operadores y asociarles precedencia y tipo, pero esa información pertenece al análisis sintáctico y queda fuera de este lexer \cite{swioperators}. Se reconoce primero el identificador completo: solo los lexemas exactos \texttt{is} y \texttt{mod} se reclasifican como operadores; \texttt{island}, \texttt{modulo} e \texttt{is2} son átomos. Entre números válidos, REAL tiene prioridad sobre ENTERO. El signo inicial siempre es un operador: \texttt{-25} produce dos tokens, mientras que \texttt{2e-3} produce un REAL.

Ambas clases de citas permiten contenido vacío y los escapes \verb|\n|, \verb|\t|, \verb|\r|, \verb|\\|, \verb|\'| y \verb|\"|. Se conserva el lexema exacto, incluidas comillas y barras, sin evaluar los escapes. Un CR o LF literal interrumpe una cita sin cierre. Los comentarios no se reconocen dentro de citas; las comillas no tienen significado especial dentro de comentarios. La expresión de bloque impide consumir un cierre intermedio, por lo que \texttt{/* a */ b /* c */} contiene dos comentarios separados por el átomo \texttt{b}.

La aceptación numérica incorpora una regla de frontera: una letra ASCII, guion bajo o un segundo punto seguido de dígito ASCII inmediatamente después del número origina un diagnóstico de número mal formado. Así, \texttt{12abc} y \texttt{3.14.15} no se fragmentan en tokens válidos. Un punto sin dígito posterior queda disponible como PUNTO: \texttt{12.} produce ENTERO y PUNTO. Esta comprobación y la recuperación de errores complementan los patrones de números válidos; no analizan la gramática de cláusulas.

\begin{table}[H]
\centering
\small
\caption{Nombres concretos de los tokens delimitadores. Cada patrón es un carácter literal.}
\label{tab:delimitadores}
\begin{tabularx}{\textwidth}{Y C{1.3cm} Y C{1.3cm}}
\toprule
\tableheader{Token} & \tableheader{Lexema} & \tableheader{Token} & \tableheader{Lexema} \\
\midrule
PARENTESIS\_IZQ & \texttt{(} & PARENTESIS\_DER & \texttt{)} \\
\rowcolor{palegray} CORCHETE\_IZQ & \texttt{[} & CORCHETE\_DER & \texttt{]} \\
LLAVE\_IZQ & \texttt{\{} & LLAVE\_DER & \texttt{\}} \\
\rowcolor{palegray} BARRA\_VERTICAL & \texttt{|} & COMA & \texttt{,} \\
PUNTO & \texttt{.} & & \\
\bottomrule
\end{tabularx}
\end{table}

Los ejemplos inválidos de la Tabla~\ref{tab:ejemplos} no pertenecen a la categoría de su fila; algunos sí son válidos como otra categoría. Por ejemplo, \texttt{Padre} es una variable y \texttt{\_\_} también. Los comentarios y espacios no emiten tokens ni atributos; los índices de átomos, variables y literales comienzan en cero dentro de sus respectivas tablas.

\small
\begin{longtable}{L{3.1cm}L{3.3cm}L{4.3cm}L{3.5cm}}
\caption{Ejemplos y atributos asociados.}\label{tab:ejemplos}\\
\toprule
\tableheader{Token} & \tableheader{Ejemplos válidos} & \tableheader{Ejemplos inválidos} & \tableheader{Atributo} \\
\midrule
\endfirsthead
\toprule
\tableheader{Token} & \tableheader{Ejemplos válidos} & \tableheader{Ejemplos inválidos} & \tableheader{Atributo} \\
\midrule
\endhead
ATOMO & \texttt{padre}, \texttt{persona\_1} & \texttt{Padre}, \texttt{1persona} & Índice en átomos \\
\rowcolor{palegray} \makecell[l]{ATOMO\_\\CITADO} & \texttt{'Juan Pérez'}, \texttt{':-'} & \texttt{'sin cierre} & Índice en átomos \\
VARIABLE & \texttt{X}, \texttt{Persona}, \texttt{\_Tmp} & \texttt{x}, \texttt{2X} & Índice en variables \\
\rowcolor{palegray} \makecell[l]{VARIABLE\_\\ANONIMA} & \texttt{\_} & \texttt{\_\_} como anónima & Sin índice \\
ENTERO & \texttt{25}, \texttt{0} & \texttt{12abc}, \texttt{2e} & Índice en literales \\
\rowcolor{palegray} REAL & \texttt{3.14}, \texttt{2e-3} & \texttt{3.14.15}, \texttt{1e+} & Índice en literales \\
CADENA & \texttt{"hola"}, \verb|"a\n b"| & \texttt{"sin cierre}, \verb|"x\q"| & Índice en literales \\
\rowcolor{palegray} \makecell[l]{OPERADOR\_\\CLAUSULA} & \texttt{:-}, \texttt{-->} & \texttt{:}, \texttt{?-} & Sin índice \\
\makecell[l]{OPERADOR\_\\CONSULTA} & \texttt{?-} & \texttt{?}, \texttt{:-} & Sin índice \\
\rowcolor{palegray} \makecell[l]{OPERADOR\_\\COMPARACION} & \verb|\==|, \texttt{=..}, \texttt{=<} & \texttt{is}, \texttt{@} & Sin índice \\
\makecell[l]{OPERADOR\_\\ARITMETICO} & \texttt{//}, \texttt{**}, \texttt{is} & \texttt{island}, \texttt{=} & Sin índice \\
\rowcolor{palegray} \makecell[l]{OPERADOR\_\\CONTROL} & \verb|\+|, \texttt{!}, \texttt{;} & \texttt{+}, \texttt{?} & Sin índice \\
DELIMITADORES & \texttt{( ) [ ] \{ \} | , .} & \texttt{@}, \texttt{\#} & Sin índice \\
\rowcolor{palegray} COMENTARIO & \texttt{\% nota}, \texttt{/* nota */} & \texttt{/* sin cierre} & Se descarta \\
ESPACIO & Espacio, TAB, CR, LF, FF, VT & Espacio Unicode U+00A0 & Se descarta \\
\bottomrule
\end{longtable}
\normalsize
